\documentclass[UTF8,fontset=none,11pt,a4paper]{ctexart}
\usepackage{ipara-handbook}
\handbooksetup{DS-10 Hash 与直接定位}{408 · 数据结构 DS-10}{Key · Collision · Load · Expected Cost}
\begin{document}
\handbooktitle{Hash 与直接定位}{用槽位和冲突管理换取接近直接访问}{408 · 数据结构 Topic DS-10}

\begin{corebox}{母问题：映射不是定位，冲突怎样成为状态的一部分？}
散列的生成链是
\[
\text{Key}\to\text{Hash Mapping}\to\text{Collision}\to\text{Resolution}\to\text{Load State}\to\text{Expected Cost}.
\]
哈希函数只给出候选槽位，不保证唯一；开放定址沿探测序列寻找空位，链地址把同槽 key 放入桶链。查找失败必须按同一探测规则走到“从未使用”或完整探测，删除不能简单清空开放定址槽位，否则会截断后续 key 的探测链。
\end{corebox}
\begin{methodbox}{本册定位与 Owner 边界}
\begin{itemize}
\item \kw{Owns：}散列函数、装填因子、开放定址/链地址、删除标记、重散列与成功/失败查找语义。
\item \kw{Uses：}Atlas Foundation 与 DS-B02 workload 比较；精确等值查询优先，不承诺范围顺序。
\item \kw{Stops：}有序索引归 DS09；具体工程哈希表的并发、缓存和安全策略只作 Extension。
\end{itemize}
\end{methodbox}
\tableofcontents\newpage

\section{状态与冲突策略}
开放定址槽位至少有 Empty、Occupied、Deleted 三态；链地址槽位可用容器表示。装填因子 $\alpha=n/m$ 越高，冲突和探测长度通常越大；超过阈值应扩容并重新插入，不能直接把旧下标复制到新表。删除标记仍占探测路径，重散列时可被清除。
\begin{examplebox}{为什么删除后不能变 Empty}
若 $h(a)=h(b)$，先插入 $a$ 再插入 $b$，$b$ 可能探测到下一个槽。删除 $a$ 后若把槽清空，查找 $b$ 看到第一个 Empty 就会错误失败；Deleted 表示“继续探测，但此处可复用”。
\end{examplebox}

\section{成本与边界}
均匀散列、合理装填因子和独立探测假设下，开放定址成功/失败查找期望接近常数；最坏可退化到 $O(m)$。哈希不保序，因此范围查询、前驱后继和有序遍历不能由平均 $O(1)$ 推出。

\section{代码与测试证据}
完整实现位于 \codepath{code/ds10_hash_table.hpp}，测试位于 \codepath{code/ds10_hash_table_test.cpp}。
\lstinputlisting[language=C++,firstline=12,lastline=108,firstnumber=12,numbers=left,caption={线性探测、删除标记与重散列}]{30_408/10_数据结构/10_Hash与直接定位/code/ds10_hash_table.hpp}
\lstinputlisting[language=C++,firstline=8,lastline=47,firstnumber=8,numbers=left,caption={冲突、删除后查找与满表边界测试}]{30_408/10_数据结构/10_Hash与直接定位/code/ds10_hash_table_test.cpp}
\begin{lstlisting}[language=bash]
clang++ -std=c++17 -Wall -Wextra -Werror -pedantic code/ds10_hash_table_test.cpp -o /tmp/ds10_test
/tmp/ds10_test
\end{lstlisting}

\section{做题控制协议}
先问查询是否精确等值，再写 $h(k)$、冲突处理、装填因子和失败终止条件。开放定址题画探测序列，删除保留 tombstone；扩容题重新计算每个 key 的槽位。若题目要求范围或排序，转交 DS09/DS11。
\begin{corebox}{压缩复原}
五问：key 映射到哪里？冲突后走哪条路径？槽位有几种状态？失败何时停止？装填变大如何修复？
\end{corebox}
\end{document}
